\section{Manage the Regime, Not Just the Peak}

A factory, a hospital, and a transport network can improve a local optimum while making their interfaces harder to manage. The next question is how to keep systems with different response times working together.

\subsection{Several Clocks in One System}

Symbols can be copied quickly. A transformer can take a long time to procure. Clinical effects may take longer to characterize. The clocks belong to particular processes. A nerve impulse is fast. Training and validating a digital system can be slow.

For a specified process $i$, a useful ratio is
\begin{equation}
 De_i=\frac{\tau_i}{T_{\mathrm{drive}}},
\end{equation}
where $\tau_i$ is a characteristic relaxation time and $T_{\mathrm{drive}}$ the driving timescale in matching units. The Deborah number comes from material behavior. Applied to organizations or biological interventions, it is a timescale analogy whose variables need operational definitions.

When the receiving process relaxes much faster than its drive, successive changes may be absorbed separately. When the drive is fast relative to relaxation, unfinished responses can overlap and history becomes important. Neither a large nor a small ratio, by itself, proves instability.

A warehouse may change routing instructions every few seconds while staffing, training, and layout evolve over months. A grid can settle electrical disturbances rapidly while adding physical capacity takes years. A drug may engage a target quickly while long-term consequences remain unknown.

\begin{table}[ht]
\centering\small
\begin{tabularx}{\textwidth}{@{}lYYY@{}}
\toprule
Layer in a project & A relatively fast process & A slower receiving process & What can accumulate\\
\midrule
Digital & Generate a design & Integrate and verify it & Unreviewed changes\\
Industrial & Change an order & Build equipment or capacity & Backlog and commitments\\
Biological & Initiate an intervention & Observe adaptation and long-term effects & Unresolved exposure and risk\\
\bottomrule
\end{tabularx}
\caption{Timescale mismatches arise within each medium as well as between media. The project must identify the actual clocks.}
\end{table}

A single clock is often the wrong design. Fast loops can handle local disturbances, slower loops can change set points, and longer processes can review goals and capacity. Separation, buffering, and clear interfaces allow different rates to coexist.

\subsection{A Region in Which Recovery Remains Possible}

An operating system rarely needs to sit at one unchanging optimum. It needs to remain in a region where service continues, disturbances are absorbed, and recovery remains available.

The fluid image suggests three qualitative regimes: orderly throughput, amplified fluctuations, and loss of essential connection. In a real project these must be identified by measurements such as queue length, recovery time, failure correlation, reserves, and continuity of critical service. Hydrodynamic laminar flow and organizational stability do not share a universal threshold.

Local improvements can move a system toward a boundary. Shorter deadlines reduce recovery time. Lower inventories reduce buffering. More frequent policy changes reduce the period over which evidence can become interpretable. The same changes can also improve a poorly managed process. Their effect depends on the full loop.

Three design moves recur.

\textbf{Set a compatible pace.} Lengthen a decision cycle where the receiving process genuinely needs time. Elsewhere, accelerate measurement and correction so the process can tolerate a faster cycle. The aim is a fit between action and response.

\textbf{Decouple the rates.} Use buffers, reserves, staged rollout, modular boundaries, and local autonomy. A buffer's required capacity depends on the integrated excess inflow over the disturbance, rather than a rule that it must last longer than every slow process. A replenishable reserve and a one-time delay have different dynamics.

\textbf{Preserve checks that produce stability.} Verification, audit, maintenance, and accountable authorization consume resources while producing evidence, early fault detection, and recoverability. They are valuable when those functions work. Delays that supply no such benefit remain candidates for removal.

An optimization that includes reliability and recovery can choose these investments. The institutional problem appears when the benefits accrue beyond the person or department paying the immediate cost. The maintenance budget is local; the avoided collapse belongs to everyone.

Someone must hold the account across the slower cycle. Equipment life exceeds a management rotation. Ecological consequences can outlast a permit. Skill formation can outlast a product generation. Records and responsibilities need enough continuity to connect the benefit claimed today with the bill that arrives later.

The chapter's examples share an accounting discipline. A body absorbs contact uncertainty. An engineered system can reduce or relocate it. An electric vehicle changes energy flows. A fleet changes asset use. Flight changes the route and creates an airspace operation. In each case, progress is established by what the complete arrangement can now do for its total cost.

Costs can genuinely fall. A complete account establishes where and by how much. The danger is an incomplete boundary that counts an interface's speed while leaving its downstream consequences unrecorded.

That danger becomes acute when symbolic action becomes cheap. A plan, credential, image, or identity can be generated faster than anyone can determine what supports it. The next chapter examines what happens when the object, the evidence, and the accountable actor cease to line up.
